\section*{Chapitre \ref{ch:m1:getting-access-futures-options} --- Accès~: contrats à terme et options cotées}

\begin{solution}{exo:m1:getting-access-futures-options:1}
Non-membre~: $10\,000 \times 1.45 = \$14\,500$. Locataire~: $10\,000 \times 0.57 + 1\,500
= \$7\,200$.
\end{solution}

\begin{solution}{exo:m1:getting-access-futures-options:2}
Location contre passage par l'extérieur~: $1\,500/(1.45 - 0.57) = 1\,705$ côtés par mois. Détention
contre location~: $(4\,000 - 1\,500)/(0.57 - 0.45) = 20\,833$ côtés par mois.
\end{solution}

\begin{solution}{exo:m1:getting-access-futures-options:3}
E-mini~: $2 \times 1.20/12.50 = 19\%$~; $2 \times 0.47/12.50 = 8\%$~; $2 \times
0.35/12.50 = 6\%$. Micro~: $2 \times 0.22/1.25 = 35\%$~; $2 \times 0.07/1.25 = 11\%$~;
$2 \times 0.04/1.25 = 6\%$.
\end{solution}

\begin{solution}{exo:m1:getting-access-futures-options:4}
8\,400 allers-retours rapportent $8\,400 \times 0.3 \times 12.50 = \$31\,500$ brut.
Non-membre~: frais $16\,800 \times 1.45 = \$24\,360$~: net \$7\,140. Locataire~: $16\,800
\times 0.57 + 1\,500 = \$11\,076$~: net \$20\,424. La même stratégie est presque
trois fois plus rentable.
\end{solution}

\begin{solution}{exo:m1:getting-access-futures-options:5}
Au rythme normal, elle termine le mois à 49\,000, soit mille de moins. Atteindre
50\,000 rapporte 5 cents sur les 50\,000 côtés~: \$2\,500, soit \$2.50 pour chacun des
1\,000 derniers côtés, plusieurs fois les frais qui leur sont associés. Elle devrait les traiter si elle
peut le faire avec une perte inférieure à ce montant.
\end{solution}

\begin{solution}{exo:m1:getting-access-futures-options:6}
Pour de grandes positions intrajournalières, la seconde formule est plus chère sur le point qui compte~:
un délai de dix minutes pour honorer un \omterm{def:m1:financing:margincall}{appel de marge} suppose de garder la trésorerie immobilisée d'avance, ou d'être
liquidée~; deux heures permettent de se financer ailleurs. L'exigence de 150~\% s'applique
aux positions de nuit, qui sont modestes. Il faut interroger les deux sur~: la manière dont l'exposition intrajournalière est
mesurée et limitée~; ce qui déclenche la liquidation et qui décide~; les intérêts
versés sur le collatéral excédentaire~; quels collatéraux sont acceptés~; et comment ils
se sont comportés envers leurs clients lors du dernier mois sous tension.
\end{solution}

\begin{solution}{exo:m1:getting-access-futures-options:7}
La voie extérieure a désormais un coût fixe de \$1\,200~: le point mort tombe à
$(1\,500 - 1\,200)/0.88 = 341$ côtés par mois. Avec une location à \$1\,100, la location
est moins chère quel que soit le volume~: son coût fixe est plus bas \emph{et} son coût variable
est plus bas, et la fonction renvoie zéro.
\end{solution}

\begin{solution}{exo:m1:getting-access-futures-options:8}
Cinq sixièmes du bénéfice relèvent d'une allocation. Il faut demander~: pendant combien de temps le \omterm{def:m1:getting-access-futures-options:incentive}{programme d'incitation} est
garanti, et la bourse peut-elle le modifier ou y mettre fin~; l'allocation est-elle par société
ou évolue-t-elle avec le volume (un capital dix fois plus élevé ne rapporte pas dix fois une
allocation fixe)~; combien de sociétés la partagent~; quelles sont les obligations et
ce qu'elles coûtent sur un marché rapide~; que reste-t-il des \$30\,000 de revenus de fourchette après la
\omterm{def:m1:what-a-trading-firm-does:adverse-selection}{sélection adverse} qu'imposent ces obligations. L'activité extensible est ici la
\$30\,000, et elle peut être négative sans le programme.
\end{solution}

\begin{solution}{pb:m1:getting-access-futures-options:1}
\textbf{1.} 12\,600 côtés~; $12\,600 \times 0.25 \times 12.50 = \$39\,375$.
\textbf{2.} $12\,600 \times 1.45 = \$18\,270$.
\textbf{3.} Net \$21\,105~; les frais représentent 46~\% du brut.
\textbf{4.} $2 \times 1.45/12.50 = 0.23$ tick par aller-retour, contre un brut de
0.50.
\textbf{5.} 1\,705 côtés par mois, 81 par jour~: la société en traite sept fois plus.
\textbf{6.} $12\,600 \times 0.57 + 1\,500 = \$8\,682$.
\textbf{7.} Net \$30\,693~: \$9\,588 de mieux par mois, soit 45~\% de plus.
\textbf{8.} $9\,588/21 = \$457$ par jour~: les frais de demande sont remboursés en 4.4
jours de séance.
\textbf{9.} Une demande comportant la déclaration des antécédents et de la situation financière, les
règles de la bourse et la compétence disciplinaire applicable au membre à titre personnel, des
restrictions sur les comptes pouvant être négociés aux \omterm{def:m1:getting-access-futures-options:membership}{tarifs membres} (le compte propre
du membre), et le lien administratif entre la location et le
compte auprès du FCM.
\textbf{10.} 20\,833 côtés par mois, soit environ 990 par jour.
\textbf{11.} 42\,000 côtés~; brut \$13\,125. Non-membre~: $42\,000 \times 0.32 =
\$13\,440$~: net $-\$315$. Locataire~: $42\,000 \times 0.12 = \$5\,040$~: net
$+\$8\,085$ (la location est déjà payée par le E-mini).
\textbf{12.} Sur les petits contrats, dont les frais atteignent un tiers de tick pour les
non-membres, les stratégies à court horizon n'existent que pour les membres. Un non-membre
doit conserver ses positions assez longtemps pour que les frais soient négligeables face au mouvement.
\textbf{13.} Rien~: la société traite 12\,600 côtés, en deçà du seuil de
20\,000. Il lui faudrait 952 côtés par jour.
\textbf{14.} Le compter pour zéro. Un revenu qui dépend d'un programme que la bourse peut
supprimer n'est qu'une prime, non une base pour embaucher ni pour prendre des engagements de capital.
\textbf{15.} Une discrimination par les prix en faveur de ceux qui apportent du volume et
de la liquidité, et un héritage de l'origine des bourses, à l'époque où elles étaient des clubs détenus par leurs membres~:
les adhésions constituaient la propriété. Le tarif non-membre est celui que paie le client
de commodité.
\textbf{16.} Les anciens traders de parquet, leurs héritiers, les sociétés ayant acheté plus de
sièges qu'elles n'utilisent pas, et les investisseurs~: un siège rapporte un revenu de location et
peut conférer d'autres droits. Les locataires sont des sociétés dont le volume justifie le tarif mais
non le prix d'achat.
\textbf{17.} La marge. À 600 côtés par jour, la location vaut \$9\,588 par
mois~; un passage de 120~\% à 200~\% de la marge de bourse peut diviser par deux la taille que
la société peut porter, ce qui divise par deux le brut de \$39\,375, ou la contraint à lever
du capital.
\textbf{18.} Un volume assez élevé pour que les commissions du FCM et les multiples de marge
coûtent plus qu'une adhésion à une chambre de compensation~: exigences en capital, \omterm{def:m1:clearing-and-settlement:waterfall}{fonds de défaillance},
contributions, personnel des opérations et exposition directe aux défaillances des autres
membres. Pour une société de deux personnes, jamais~; pour une grande société pour compte propre, une
étape courante.
\textbf{19.} \textbf{1\,705 côtés par mois.}
\textbf{20.} Un prix plus bas par transaction, en contrepartie d'un coût fixe et de
l'acceptation de la compétence de la bourse.
\end{solution}

\begin{solution}{iq:m1:getting-access-futures-options:1}
Un \omterm{def:m1:getting-access-futures-options:fcm}{futures commission merchant} (FCM) accepte les ordres des clients et leurs fonds. Pour
une société de trading, c'est le \omterm{def:m1:getting-access-futures-options:fcm}{membre compensateur} qui porte ses positions à la
chambre de compensation, collecte et verse la marge, garantit ses transactions, fournit
ou accepte des exécutions par \omterm{def:m1:getting-access-equities:brokers}{give-up}, et fixe les limites réelles de la société~: le multiple
de marge, l'exposition intrajournalière, les droits de liquidation.
\iqlookfor{le FCM comme fournisseur de crédit, et non comme routeur d'ordres.}
\end{solution}

\begin{solution}{iq:m1:getting-access-futures-options:2}
Les bourses appartenaient à leurs membres~; les \omterm{def:m1:getting-access-futures-options:membership}{tarifs membres} subsistent sous la forme d'une remise sur volume
attachée à un siège négociable ou louable. L'écart est important~: pour le principal
contrat sur indice d'actions, un non-membre paie environ trois fois les frais de bourse d'un propriétaire
par côté, et l'écart est plus grand encore lorsqu'on le rapporte au tick des
plus petits contrats.
\iqlookfor{un ordre de grandeur de l'écart, et la location comme voie d'accès.}
\end{solution}

\begin{solution}{iq:m1:getting-access-futures-options:3}
Frais en ticks, au tarif que la société paiera réellement~: pour un non-membre, un
aller-retour sur \omterm{def:m1:index-and-single-stock-futures-worldwide:point}{contrat micro} coûte à lui seul un tiers de tick en frais de bourse, auxquels
s'ajoute la commission. Viennent ensuite la position dans la file (un scalping d'un tick suppose des exécutions en tête de
file, ce que le backtest supposait probablement), la latence, et la question de savoir si
les exécutions supposées sont réalisables au meilleur cours, tout simplement. La plupart de ces stratégies échouent à
le premier contrôle.
\iqlookfor{les frais en fraction du tick, selon le statut.}
\end{solution}

\begin{solution}{iq:m1:getting-access-futures-options:4}
La location se justifie lorsque le volume mensuel multiplié par l'économie par côté dépasse le loyer,
soit généralement quelques dizaines d'allers-retours par jour sur les principaux contrats. L'achat se justifie lorsque
l'économie du tarif propriétaire sur le tarif locataire, multipliée par le volume, dépasse le coût du
capital immobilisé, net du loyer versé, et que la société compte rester dans le
produit pendant des années~: les adhésions ont un prix qui fluctue.
\iqlookfor{les deux seuils de rentabilité et le capital à risque dans un achat.}
\end{solution}

\begin{solution}{iq:m1:getting-access-futures-options:5}
Négociables~: la commission par côté et son minimum~; le multiple de marge et
les limites intrajournalières~; les intérêts sur les liquidités excédentaires et le collatéral accepté~;
les frais de \omterm{def:m1:getting-access-equities:brokers}{give-up} et de plateforme~; les délais de préavis. Non négociables~: les frais de bourse et
réglementaires, la marge de la chambre de compensation, les règles de ségrégation. La société
apporte du volume, des soldes, un profil de risque propre et un autre \omterm{def:m1:getting-access-futures-options:fcm}{futures commission merchant} (FCM).
\iqlookfor{les conditions de marge avant la commission.}
\end{solution}

\begin{solution}{iq:m1:getting-access-futures-options:6}
Coûts fixes et échelle. Une société de futures n'a besoin que d'une adhésion, d'un flux de
données et d'un carnet d'ordres. Un \omterm{def:m1:options-market-structure:mm}{teneur de marché d'options} a besoin de désignations et d'autorisations sur une douzaine ou
davantage de bourses pour une même série~; de la capacité à coter des centaines de milliers
de séries et à absorber un flux de plusieurs centaines de milliards de messages, des modèles
pour toute la surface, du capital destiné à l'inventaire, qui ne peut être soldé à la
clôture, et assez de diversité de flux pour que l'inventaire se compense. Les droits de données
récompensent les acteurs en place. L'échelle minimale efficace est élevée, si bien que le secteur est
concentré.
\iqlookfor{l'inventaire et la diversité des flux, et pas seulement la technologie.}
\end{solution}
